% !TEX root = ../main.tex
% ======================================================================
% Section 7 – Valorisation et diffusion
% Responsable : Stive Roussel Tankio Djiokap
% Grille : « Valorisation & diffusion » – 8 pts (section prioritaire).
% Longueur indicative : ~2 à 2,5 pages.
% ======================================================================
\subsection{Valorisation}
\label{sec:valorisation-modele}
% TODO : plateforme publique de recherche ET/OU licence aux fabricants d'appareils médicaux ?
% TODO : analyser la tension entre ces deux voies.
% TODO : partage des bénéfices avec les communautés participantes.
\arediger{valorisation}

\subsection{Diffusion}
\label{sec:diffusion}

\subsubsection{Comités d'éthique de la recherche et instances de gouvernance autochtones}
\label{sec:diffusion-cer}
% TODO : public 1 — les comités d'éthique de la recherche (\acs{CER}) et les instances de gouvernance autochtones (\eptc{}, chapitre 9).
\arediger{diffusion auprès des CER et des instances autochtones}

\subsubsection{Patients et vulgarisation}
\label{sec:diffusion-patients}
% TODO : public 2 — les patients, avec un volet vulgarisation.
% TODO : comment les rassurer sur la sécurité et la pertinence du projet ?
\arediger{diffusion auprès des patients}
